\textbf{Phần 1. Trắc nghiệm lựa chọn}
\setcounterpageref{ex}{0}

\begin{ex} %[4D8N47]
Mua 5 quyển vở hết $40\,000$ đồng. Mua 1 quyển vở như thế hết số tiền là
\choice
{$7000$ đồng}
{$9000$ đồng}
{$10\,000$ đồng}
{\True $8000$ đồng}
\loigiai{
$40\,000 : 5 = 8000$ đồng.\\
\textbf{Đáp án của câu hỏi là D}
}
\end{ex}

\begin{ex} %[4D8N47]
Có $36$ chiếc bánh xếp đều vào 6 hộp. Hỏi 1 hộp có bao nhiêu chiếc bánh?
\choice
{$5$ chiếc}
{$7$ chiếc}
{\True $6$ chiếc}
{$8$ chiếc}
\loigiai{
$36 : 6 = 6$ chiếc.\\
\textbf{Đáp án của câu hỏi là C}
}
\end{ex}

\begin{ex} %[4D8H47]
Mua 4 chiếc bút chì hết $24\,000$ đồng. Mua 7 chiếc bút chì như thế hết số tiền là
\choice
{$38\,000$ đồng}
{\True $42\,000$ đồng}
{$45\,000$ đồng}
{$48\,000$ đồng}
\loigiai{
Giá 1 chiếc: $24\,000 : 4 = 6000$ đồng.\\
7 chiếc: $6000 \times 7 = 42\,000$ đồng.\\
\textbf{Đáp án của câu hỏi là B}
}
\end{ex}

\begin{ex} %[4D8H47]
Có $48\text{ kg}$ đường đóng đều vào 8 túi. Hỏi 5 túi như thế có bao nhiêu ki-lô-gam đường?
\choice
{$25\text{ kg}$}
{$35\text{ kg}$}
{$40\text{ kg}$}
{\True $30\text{ kg}$}
\loigiai{
1 túi: $48 : 8 = 6\text{ kg}$.\\
5 túi: $6 \times 5 = 30\text{ kg}$.\\
\textbf{Đáp án của câu hỏi là D}
}
\end{ex}

\begin{ex} %[4D8H47]
Một người làm việc trong 3 ngày được trả $750\,000$ đồng. Nếu người đó làm việc trong 5 ngày thì được trả số tiền là
\choice
{\True $1\,250\,000$ đồng}
{$1\,200\,000$ đồng}
{$1\,300\,000$ đồng}
{$1\,500\,000$ đồng}
\loigiai{
1 ngày: $750\,000 : 3 = 250\,000$ đồng.\\
5 ngày: $250\,000 \times 5 = 1\,250\,000$ đồng.\\
\textbf{Đáp án của câu hỏi là A}
}
\end{ex}

\begin{ex} %[4D8H47]
May 6 bộ quần áo như nhau hết $18\text{ m}$ vải. May 10 bộ quần áo như thế hết số mét vải là
\choice
{$25\text{ m}$}
{$35\text{ m}$}
{\True $30\text{ m}$}
{$40\text{ m}$}
\loigiai{
1 bộ: $18 : 6 = 3\text{ m}$.\\
10 bộ: $3 \times 10 = 30\text{ m}$.\\
\textbf{Đáp án của câu hỏi là C}
}
\end{ex}

\begin{ex} %[4D8V47]
Một ô tô đi quãng đường $120\text{ km}$ hết $3$ giờ. Với vận tốc đó, ô tô đi quãng đường $200\text{ km}$ hết số giờ là
\choice
{$4$ giờ}
{$6$ giờ}
{$7$ giờ}
{\True $5$ giờ}
\loigiai{
Mỗi giờ ô tô đi: $120 : 3 = 40\text{ km}$.\\
Thời gian đi $200\text{ km}$: $200 : 40 = 5$ giờ.\\
\textbf{Đáp án của câu hỏi là D}
}
\end{ex}

\begin{ex} %[4D8V47]
Có $72\text{ l}$ dầu đóng đều vào 9 can. Người ta lấy đi 4 can dầu đó. Số lít dầu còn lại là
\choice
{\True $40\text{ l}$}
{$32\text{ l}$}
{$48\text{ l}$}
{$36\text{ l}$}
\loigiai{
1 can: $72 : 9 = 8\text{ l}$.\\
Số can còn lại: $9 - 4 = 5$ can.\\
Số dầu còn lại: $8 \times 5 = 40\text{ l}$.\\
\textbf{Đáp án của câu hỏi là A}
}
\end{ex}

\begin{ex} %[4D8V47]
Mua 8 hộp sữa chua hết $64\,000$ đồng. Bạn Lan có $40\,000$ đồng thì mua được nhiều nhất bao nhiêu hộp sữa chua như thế?
\choice
{$4$ hộp}
{$6$ hộp}
{\True $5$ hộp}
{$7$ hộp}
\loigiai{
Giá 1 hộp: $64\,000 : 8 = 8000$ đồng.\\
Số hộp: $40\,000 : 8000 = 5$ hộp.\\
\textbf{Đáp án của câu hỏi là C}
}
\end{ex}

\begin{ex} %[4D8V47]
Một đội thợ xây lát $45\text{ m}^2$ sân hết $1125$ viên gạch. Để lát $60\text{ m}^2$ sân như thế cần số viên gạch là
\choice
{$1400$ viên}
{\True $1500$ viên}
{$1600$ viên}
{$1800$ viên}
\loigiai{
$1\text{ m}^2$ cần: $1125 : 45 = 25$ viên.\\
$60\text{ m}^2$ cần: $25 \times 60 = 1500$ viên.\\
\textbf{Đáp án của câu hỏi là B}
}
\end{ex}

\begin{ex} %[4D8V47]
Một vòi nước chảy trong 4 phút được $60\text{ l}$ nước. Để chảy được $150\text{ l}$ nước vào bể thì cần thời gian là
\choice
{\True $10$ phút}
{$8$ phút}
{$12$ phút}
{$15$ phút}
\loigiai{
1 phút chảy: $60 : 4 = 15\text{ l}$.\\
Thời gian: $150 : 15 = 10$ phút.\\
\textbf{Đáp án của câu hỏi là A}
}
\end{ex}

\begin{ex} %[4D8C47]
Một nhóm thợ gồm 6 người làm trong 4 ngày được 120 sản phẩm. Hỏi 8 người thợ làm việc với năng suất như thế trong 5 ngày thì được bao nhiêu sản phẩm?
\choice
{$180$ sản phẩm}
{\True $200$ sản phẩm}
{$220$ sản phẩm}
{$240$ sản phẩm}
\loigiai{
1 người làm trong 1 ngày: $120 : (6 \times 4) = 5$ sản phẩm.\\
8 người làm trong 5 ngày: $8 \times 5 \times 5 = 200$ sản phẩm.\\
\textbf{Đáp án của câu hỏi là B}
}
\end{ex}

\textbf{Phần 2. Trắc nghiệm đúng sai}
\setcounterpageref{ex}{0}

\begin{ex} %[4D8H47]
Mua 6 quyển sách hết $120\,000$ đồng. Xét tính đúng sai:
\choiceTF
{\True Giá 1 quyển sách là $20\,000$ đồng.}
{\True Mua 10 quyển sách như thế hết $200\,000$ đồng.}
{Với $300\,000$ đồng mua được 12 quyển sách.}
{\True Với $300\,000$ đồng mua được 15 quyển sách.}
\loigiai{
\begin{itemize}
    \item 1 quyển: $120\,000 : 6 = 20\,000$ đồng. Đúng.
    \item 10 quyển: $20\,000 \times 10 = 200\,000$ đồng. Đúng.
    \item Với 300k mua được $300\,000 : 20\,000 = 15$ quyển $\neq 12$ quyển. Ý 3 Sai, Ý 4 Đúng.
\end{itemize}
\begin{itemize}
    \item a) \textbf{Đúng}.
    \item b) \textbf{Đúng}.
    \item c) \textbf{Sai}.
    \item d) \textbf{Đúng}.
\end{itemize}
}
\end{ex}

\begin{ex} %[4D8V47]
Một ô tô tiêu thụ hết 12 lít xăng khi đi quãng đường $150\text{ km}$. Xét tính đúng sai:
\choiceTF
{\True Đi quãng đường $50\text{ km}$ ô tô tiêu thụ hết 4 lít xăng.}
{\True Đi quãng đường $300\text{ km}$ ô tô tiêu thụ hết 24 lít xăng.}
{Với 6 lít xăng ô tô đi được quãng đường $80\text{ km}$.}
{\True Với 6 lít xăng ô tô đi được quãng đường $75\text{ km}$.}
\loigiai{
\begin{itemize}
    \item $50\text{ km}$ cần: $12 : 3 = 4$ lít xăng. Đúng.
    \item $300\text{ km}$ gấp đôi $150\text{ km} \Rightarrow 12 \times 2 = 24$ lít. Đúng.
    \item 6 lít xăng $= 12 : 2 \Rightarrow 150 : 2 = 75\text{ km} \neq 80\text{ km}$. Ý 3 Sai, Ý 4 Đúng.
\end{itemize}
\begin{itemize}
    \item a) \textbf{Đúng}.
    \item b) \textbf{Đúng}.
    \item c) \textbf{Sai}.
    \item d) \textbf{Đúng}.
\end{itemize}
}
\end{ex}

\textbf{Phần 3. Trả lời ngắn}
\setcounterpageref{ex}{0}

\begin{ex} %[4D8H47]
Mua 5 cái bánh hết $35\,000$ đồng. Tính tiền mua 8 cái bánh như thế theo nghìn đồng (chỉ ghi số nghìn đồng).
\shortans{56}
\loigiai{
\textbf{Đáp số:} $56$\\
$35 : 5 \times 8 = 56$ nghìn đồng.
}
\end{ex}

\begin{ex} %[4D8H47]
Có $42\text{ kg}$ gạo chia đều vào 7 túi. Hỏi 9 túi như thế chứa bao nhiêu ki-lô-gam gạo?
\shortans{54}
\loigiai{
\textbf{Đáp số:} $54$\\
$42 : 7 \times 9 = 54\text{ kg}$.
}
\end{ex}

\begin{ex} %[4D8V47]
Một người thợ may trong 4 ngày may được 16 bộ quần áo. Hỏi trong 9 ngày người đó may được bao nhiêu bộ quần áo?
\shortans{36}
\loigiai{
\textbf{Đáp số:} $36$\\
$16 : 4 \times 9 = 36$ bộ.
}
\end{ex}

\begin{ex} %[4D8V47]
Một vòi nước chảy trong 5 phút được $100\text{ l}$ nước. Hỏi vòi đó chảy trong 12 phút được bao nhiêu lít nước?
\shortans{240}
\loigiai{
\textbf{Đáp số:} $240$\\
$100 : 5 \times 12 = 240$ lít.
}
\end{ex}

\begin{ex} %[4D8V47]
Có $60\text{ l}$ mật ong đóng đều vào 12 can. Hỏi có $45\text{ l}$ mật ong thì đóng được bao nhiêu can như thế?
\shortans{9}
\loigiai{
\textbf{Đáp số:} $9$\\
1 can: $60 : 12 = 5\text{ l} \Rightarrow 45 : 5 = 9$ can.
}
\end{ex}

\begin{ex} %[4D8C47]
Một xí nghiệp có 15 công nhân may được 300 sản phẩm trong một ngày. Nếu tăng thêm 5 công nhân nữa thì mỗi ngày xí nghiệp may được bao nhiêu sản phẩm?
\shortans{400}
\loigiai{
\textbf{Đáp số:} $400$\\
1 công nhân: $300 : 15 = 20$ sản phẩm.\\
20 công nhân: $20 \times 20 = 400$ sản phẩm.
}
\end{ex}
